شهد \textbf{قطاع الخدمات المالية} أفضل مستوى أداء؛ حيث سجل قيما تجاوزت \textbf{المستوى المحايد} 20نقطة. ويرجع ذلك إلى ارتفاع معدلات التداول نتيجة استمرار تطبيق البنك المركزي لسياسات التيسير النقدي، بالإضافة إلى إلغاء ضريبة الأرباح الرأسمالية على تعاملات البورصة بهدف تشجيع الاستثمار في سوق المال. كما ساهم في ارتفاع معدلات التداول الاستقرار الاقتصادي الكلي للمؤشرات وتعزيز ثقة المستثمرين في الاقتصاد المصري. إلى جانب تحسن التصنيف الائتماني لمصر من قِبل مؤسسات التصنيف الدولية.
